%قبل از شروع مسئله ترجیحاً چیزی ننویسید. اما اگر خواستید بنویسید، در قالب لیستی (فقرات و شمارش) بنویسید.
%برای اضافه کردن نمره، بعدی از شروع‌مسئله، [x] را اضافه کنید و به جای x  نمره مورد نظر را قرار دهید. برای زیرمسئله‌ها هم به همین صورت است.
%\شروع‌مسئله[10]
\شروع‌مسئله [15]
% صورت اصلی مسئله را در این‌جا بنویسید.
رشته‌ای به طول $n$ داریم. الگوریتمی ارائه دهید که تعداد زیررشته‌های متقارن این رشته را در زمان
$\mathcal{O}(n \log n)$
	محاسبه کند. برای این سوال یک تقریب مناسب از جواب نیز کفایت می‌کند، به عبارت دیگر در صورت پیاده‌سازی یک تابع درهم‌سازی مناسب نیازی نیست حالاتی که  \متن‌لاتین{hash} برابر است اما رشته‌ها برابر نیستند را بررسی کنید. همچنین
	منظور از رشتۀ متقارن رشته‌ای است که از دو طرف به یک شکل خوانده شود مانند \متن‌لاتین{abccba} و یا \متن‌لاتین{aca}.  

%اگر مسئله چند قسمت داشت به این صورت عمل کنید:
%\شروع{زیرمسائل}
%\زیرمسئله[0] % مسئلۀ ۱آ
%\زیرمسئله % مسئلۀ ۱ب
%\زیرمسئله % مسئلۀ ۱پ
%\زیرمسئله % مسئلۀ ۱ت
%\پایان{زیرمسائل}

\solution{
رشته  مسئله را 
$s$
می‌نامیم.
برای این سؤال ابتدا درهم‌سازی تمامی زیررشته‌های به فرم
$s[0, t]$
را محاسبه می‌کنیم. برای این کار درهم‌سازی
را به صورت زیر در نظر می‌گیریم (برای $m$ و $k$ مناسب)
\begin{equation*}
	h(s[0, t]) = \sum_{i = 0}^{t} s[i]k^{i} \mod m 
\end{equation*}
دقت کنید  محاسبه تمامی این درهم‌سازی‌ها در زمان 
$\mathcal{O}(n)$
انجام می‌شود. همین کار را برای وارونه رشته $s$ (که آن را $s'$ می‌نامیم) نیز انجام می‌دهیم. توجه کنید با این کار درهم‌سازی یک زیررشتهٔ خاص در زمان 
 $\mathcal{O}(1)$
 پیدا می‌شود، زیرا
 \begin{equation*}
 	h(s[i, j]) = h(s[0, j]) - h(s[0, i-1])
 \end{equation*}
 پس درهم‌سازی زیررشتهٔ
 $s[i, j]$
 در زمان
 $\mathcal{O}(1)$
محاسبه می‌شود. به طور مشابه درهم‌سازی وارون آن نیز در 
  $\mathcal{O}(1)$ 
  محاسبه می‌شود، پس می‌توان متقارن بودن آن را در 
  $\mathcal{O}(1)$
  بررسی کرد. پس تا اینجا نشان دادیم که بررسی متقارن بودن هر زیررشته
  در 
  $\mathcal{O}(1)$
 انجام می‌شود.

  ابتدا تعداد زیررشته‌های متقارن با طول فرد را پیدا می‌کنیم. توجه کنید هر یک از این زیررشته‌ها یک حرف وسط دارد. اندیس این حرف وسط را $k$ می‌گیریم. حال توجه کنید که اگر 
  $s[k - i, k +i]$
  متقارن باشد؛ تمام زیررشته‌های کوچک‌تر با حرف وسط
  $k$
  نیز متقارن هستند؛ به عبارت دیگر 
 $s[k - i + 1, k +i - 1]$
 و 
 \dots
 و
 $s[k - 1, k + 1]$
 نیز متقارن هستند
 (می‌توان خود حرف را نیز به عنوان زیررشته متقارن حساب کرد). پس کافی است برای هر حرف وسطی، بزرگترین زیررشته متقارن به مرکز آن حرف را پیدا کنیم. بنابر توضیحات قبلی برای این کار می‌توانیم از جست‌و‌جوی دودویی استفاده کنیم. با استفاده از جست‌و‌جوی دودویی کافی است برای هر حرف وسط تنها متقارن بودن
$\mathcal{O}(\log n)$
زیررشته را بررسی کنیم، که هرکدام در زمان 
$\mathcal{O}(1)$
انجام می‌شود. پس تعداد زیررشته‌های متقارن به مرکز یک حرف خاص در زمان 
$\mathcal{O}(\log n)$
محاسبه می‌شود. از آن‌جا که در کل 
$n$
حرف داریم، تعداد کل زیررشته‌های متقارن فرد در زمان 
$\mathcal{O}(n\log n)$
قابل محاسبه است.
\\
برای راحتی، یک تعریف می‌کنیم. به فاصله بین دو حرف از رشته یک  دو-فاصله می‌گوییم. در حقیقت مسئله برای زیررشته‌های به طول زوج فرق خاصی با زیررشته‌های به طول فرد ندارد. کافی است توجه کنیم که زیررشته‌های زوج حرف وسط ندارند؛ اما دو-فاصله وسط دارند. مشابه جست‌وجوی دودویی را می‌توان در مورد دو-فاصله‌ها به کار برد.
با استفاده از جست‌و‌جوی دودویی کافی است برای هر دو-فاصله وسط تنها متقارن بودن
$\mathcal{O}(\log n)$
زیررشته را بررسی کنیم، که هرکدام در زمان 
$\mathcal{O}(1)$
انجام می‌شود. پس تعداد زیررشته‌های متقارن به مرکز یک دو-فاصله خاص در زمان 
$\mathcal{O}(\log n)$
محاسبه می‌شود. از آن‌جا که در کل 
$n - 1$
دو-فاصله داریم، تعداد کل زیررشته‌های متقارن زوج نیز در زمان 
$\mathcal{O}(n\log n)$
محاسبه می‌شود. با جمع زدن تمامی زمان‌های فوق می‌توان دید تعداد کل زیررشته‌های متقارن در زمان 
$\mathcal{O}(n\log n)$
محاسبه می‌شود.

}




%%%%%%%%%%%%%%%%%%%%%%%%%%%%%%%%%%%%%%%%%%%%%%%%%%%%%
% برخی دستورات مفید برای تمرین‌ها. %
%%%%%%%%%%%%%%%%%%%%%%%%%%%%%%%%%%%%%%%%%%%%%%%%%%%%%

% فرمول‌های نمونه:
%$f(x) = \Theta(x)$
%$T(x, y) \leq \log(x) + 2^y + \binom{2n}{n}$
% {\tt code\_function}


% لیست موارد بدون شماره‌گذاری:
%\شروع{فقرات}
%    \فقره مورد اول 
%        \شروع{فقرات}
%            \فقره مورد اول داخلی 
%                \شروع{فقرات}
%                    \فقره اولین مورد خیلی داخلی 
%                    \فقره دومین مورد خیلی داخلی 
%                \پایان{فقرات}
%            \فقره مورد دوم داخلی 
%        \پایان{فقرات}
%    \فقره مورد دوم 
%\پایان{فقرات}

% لیست موارد شماره‌گذاری شده:
%\شروع{شمارش}
%    \فقره یک مورد 
%    \فقره یک مورد دیگر 
%    \فقره باز هم یک مورد دیگر
%\پایان{شمارش}

% برای نوشتن یک فرمول چند خطی که یک شماره بخوره. برای حالت‌های دیگر به سایت زیر مراجعه کنید:
% https://www.overleaf.com/learn/latex/Aligning_equations_with_amsmath
%\begin{align} 
%    \begin{split}
%        (x+y)^3 &= (x+y)^2(x+y) \\
%                &= (x^2+2xy+y^2)(x+y) \\
%                &= (x^3+2x^2y+xy^2) + (x^2y+2xy^2+y^3) \\
%                &= x^3+3x^2y+3xy^2+y^3
%    \end{split}                                 
%\end{align}

% نوشتن ماتریس و آرایه دوبعدی
%\begin{align}
%    A = 
%    \begin{bmatrix}
%        A_{11} & A_{21} \\
%        A_{21} & A_{22}
%    \end{bmatrix}
%\end{align}

% اضافه کردن تصویر با شرح. برای اضافه کردن تصویر بدون شرح، فقط دستور درج‌تصویر را استفاده کنید.
%\شروع{شکل}[h!]
%\تنظیم‌ازوسط
%  \درج‌تصویر[پهنا=0.5\پهنای‌متن]{img.jpg}
%\شرح{شرح تصویر}
%\برچسب{تصویر:۱}
%\پایان{شکل}
%\پایان{زیرمسائل}

%کد زیر را ببینید.
%%\lstinputlisting{source_filename.py}
%\begin{latin}
%\begin{lstlisting}
%# Import the modules
%import sys
%import random
%
%ans = True
%
%while ans:
%question = raw_input("Ask the magic 8 ball a question: (press enter to quit) ")
%
%answers = random.randint(1,8)
%
%if question == "":
%sys.exit()
%
%elif answers == 1:
%print "It is certain"
%
%elif answers == 2:
%print "Outlook good"
%
%elif answers == 3:
%print "You may rely on it"
%
%elif answers == 4:
%print "Ask again later"
%
%elif answers == 5:
%print "Concentrate and ask again"
%
%elif answers == 6:
%print "Reply hazy, try again"
%
%elif answers == 7:
%print "My reply is no"
%
%elif answers == 8:
%print "My sources say no"
%\end{lstlisting}
%\end{latin}
